\chapter{Controlled Ricci flows with surgery}
\label{chap:controlled-ricci-flows}

 We do not
wish to consider all Ricci flows with surgery. Rather we shall
concentrate on $3$-dimensional flows (that is to say $4$-dimensional
space-times) whose singularities are closely controlled both
topologically and geometrically. We introduce the hypotheses that we
require these evolutions to satisfy. Then main result, which is
stated in this chapter and proved in the next two, is that these
controlled $3$-dimensional Ricci flows with surgery always exist for
all time with any compact $3$-manifold as initial metric.

\subsection{Normalized initial conditions}

Consider a compact connected Riemannian $3$-manifold $(M,g(0))$
satisfying
\begin{enumerate}
\item[(1)] $|\Rm(x,0)|\le 1$ for all $x\in M$ and
\item[(2)] for every $x\in M$ we have $\Vol\,B(x,0,1)\ge \omega/2$ where
$\omega$ is the volume of the unit ball in $\Ar^3$.
\end{enumerate}

Under these conditions we say that $(M,g(0))$ is {\em normalized}.
Also, if $(M,g(0))$ is the initial manifold of a Ricci flow with
surgery then we say that it is a {\em normalized initial metric}. Of
course, given any compact Riemannian $3$-manifold $(M,g(0))$ there
is a positive constant $Q<\infty$ such that $(M,Qg(0))$ is
normalized.




Starting with a normalized initial metric implies that the flow exists
and has uniformly bounded curvature for a fixed amount of time. This is the
content of the following claim which is an immediate corollary of
Theorem~\ref{compuniq}, Proposition~\ref{patch}, Theorem~\ref{shi},
and  Proposition~\ref{kappa0r0t0}.


\begin{claim}[Normalized initial data give short-time existence with curvature and volume control]
\label{claim:normalized-initial-existence}\label{newkappa0r0t0}
\uses{thm:local-existence-uniqueness,prop:smooth-patching,thm:shi-derivative-estimates,prop:normalized-flow-short-time-bounds}
 There is $\kappa_0$ such that the following holds. Let $(M,g(0))$ be a normalized initial
metric. Then the solution to the Ricci flow equation with these
initial conditions exists for $t\in [0,2^{-4}]$, and  $|R(x,t)|\le
2$ for all $x\in M$ and all $t\in [0,2^{-4}]$. Furthermore, for any
$t\in [0,2^{-4}]$ and any $x\in M$ and any $r\le \epsilon$ we have
$\Vol\,B(x,t,r)\ge \kappa_0r^3$.
\end{claim}


\section{Gluing together evolving necks}



\begin{proposition}[Gluing an evolving neck to a strong neck produces a longer strong neck]
\label{prop:neck-gluing}
\uses{def:strong-canonical-neighborhood-assumption,def:evolving-epsilon-neck}\label{neckglue}
There is $0<\beta<1/2$ such that the following holds for any
$\epsilon<1$. Let $(N\times [-t_0,0],g_1(t))$ be an evolving
$\beta\epsilon$-neck centered at $x$ with $R(x,0)=1$. Let $
(N'\times (-t_1,-t_0],g_2(t))$ be a strong $\beta\epsilon/2$-neck.
Suppose we have an isometric embedding of $N\times \{-t_0\}$ with
$N'\times \{-t_0\}$ and the strong $\beta\epsilon/2$-neck structure on $N'\times
(-t_1,-t_0]$ is centered at the image of $(x,-t_0]$. Then the union
$$N\times [-t_0,0]\cup N'\times (-t_1,-t_0]$$ with the induced
one-parameter family of metrics contains a strong $\epsilon$-neck
centered at $(x,0)$.
\end{proposition}


\begin{proof}
Suppose that the result does not hold. Take a sequence of $\beta_n$
tending to zero and counterexamples $(N_n\times
[-t_{0,n},0],g_{1,n}(t));\ (N'_n\times
(-t_{1,n},-t_{0,n}],g_{2,n}(t))$. Pass to a subsequence so that the
$t_{0,n}$ tend to a limit $t_{0,\infty}\ge 0$. Since $\beta_n$ tends
to zero, we can take a smooth limit of a subsequence and this limit
is an evolving cylinder $(S^2\times \Ar, h(t)\times ds^2)$, where
$h(t)$ is the round metric of scalar curvature $1/(1-t)$ defined for
some amount of backward time. Notice that, for all $\beta$
sufficiently small, on a $\beta\epsilon$-neck the derivative of the
scalar curvature is positive. Thus, $R_{g_{1,n}}(x,-t_{0,n})<1$.
Since we have a strong neck structure on $N'_n$ centered at
$(x,-t_{0,n})$, this implies that $t_{1,n}>1$ so that the limit is
defined for at least time $t\in [0, 1+t_{0,\infty})$. If
$t_{0,\infty}>0$, then, restricting to the appropriate subset of
this limit, a subset with compact closure in space-time, it follows
immediately that for all $n$ sufficiently large there is a strong
$\epsilon$-neck centered at $(x,0)$. This contradicts the assumption
that we began with a sequence of counterexamples to the
proposition.

Let us consider the case when $t_{0,\infty}=0$. In this case the
smooth limit is an evolving round cylinder defined for time
$(-1,0]$.
 Since $t_{1,n}>1$ we
see that for any $A<\infty$ for all $n$ sufficiently large the ball
$B(x_n,0,A)$ has compact closure in every time-slice and there are
uniform bounds to the curvature on $B(x_n,0,A)\times (-1,0]$. This
means that the limit is uniform for time $(-1,0]$ on all these
balls. Thus, once again for all $n$ sufficiently large we see that
$(x,0)$ is the center of a strong $\epsilon$-neck in the union. In
either case we have obtained a contradiction, and hence we have
proved the result. 
\end{proof}








\subsection{First assumptions}\label{firstassumpt}

\noindent{\bf Choice of $C$ and $\epsilon$:} The first thing we need
to do is fix for the rest of the argument $C<\infty$ and
$\epsilon>0$. We do this in the following way. We fix $0<\epsilon\le
\mathrm{min}(1/200,\left(\sqrt{D}(A_0+5)\right)^{-1},\bar\epsilon_1/2,\bar
\epsilon'/2,\epsilon_0)$ where  $\bar\epsilon_1$ is the constant from
Proposition~\ref{narrows}, $\bar\epsilon'$ is the constant from
Theorem~\ref{kappasummary}, $\epsilon_0$ is the constant from Section~\ref{10.1},
and $A_0$ and $D$ are the constants from
Lemma~\ref{A_0}.
 We fix $\beta<1/2$, the constant from
Proposition~\ref{neckglue}.  Then we let $C$ be the maximum of the
constant $C(\epsilon)$ as in Corollary~\ref{kappacannbhd} and
$C'(\beta\epsilon/3)+1$ as in Theorem~\ref{stdsolncannbhd}.

For all such $\epsilon$, Theorem~\ref{bcbd} holds for $\epsilon$ and
Proposition~\ref{narrows}, Proposition~\ref{canonvary} and
Corollaries~\ref{kappacannbhd} and~\ref{limitcannbhd} and Theorems~\ref{smlmtflow} and~\ref{kaplimit}
 hold for $2\epsilon$. Also, all the topological results of the Appendix hold
for $2\epsilon$ and $\alpha=10^{-2}$.


Now let us turn to the assumptions we shall make on the Ricci flows
with surgery that we shall consider. Let ${\mathcal M}$ be a
space-time. Our first  set of assumptions are basically topological
in nature. They are:

\bigskip

\noindent {\bf Assumption (1). Compactness and dimension:} {\em Each
time-slice $M_t$ of space-time is a compact $3$-manifold
containing no embedded $\Ar P^2$ with trivial normal bundle.}


\bigskip


\noindent{\bf Assumption (2). Discrete singularities:} {\em The set of
singular times is a discrete subset of  $\Ar$.}


\bigskip



\noindent {\bf Assumption (3). Normalized initial conditions:} {\em
$0$ is the initial time of the Ricci flow with surgery and the
initial metric $(M_0,G(0))$ is normalized.}

\bigskip


 It follows
from Assumption (2) that for any time $t$ in the time-interval of
definition of a Ricci flow with surgery, with $t$ being distinct
from the initial and final times (if these exist), for all
$\delta>0$ sufficiently small, the only possible singular time in
$[t-\delta,t+\delta]$ is $t$. Suppose that $t$ is a singular time.
The singular locus at time $t$ is a closed, smooth subsurface
$\Sigma_{t}\subset M_{t}$. From the local
model, near every point of $x\in \Sigma_{t}$ we see that this
surface separates $M_{t}$ into two pieces:
$$M_{t}=C_{t}\cup_{\Sigma_{t}}E_{t},$$
where $E_{t}$ is the exposed region at time $t$ and $C_{t}$ is the
complement of the interior of $E_{t}$ in $M_{t}$. We call $C_t$ the
{\em continuing region}.  $C_t\subset M_t$ is the maximal subset of
$M_t$ for which there is $\delta>0$ and an embedding $C_t\times
(t-\delta,t]\to {\mathcal M}$ compatible with time and the vector
field.

\noindent {\bf Assumption (4). Topology of the exposed regions:} {\em
At all singular times $t$  we require that $E_{t}$ be a finite
disjoint union of $3$-balls. In particular, $\Sigma_{t}$ is a
finite disjoint union of $2$-spheres.}

\bigskip

The next assumptions are geometric in nature.
Suppose that $t$ is a surgery time. Let ${\mathcal M}_{(-\infty,t)}$ be ${\bf t}^{-1}((-\infty,t))$
and let $(\widehat {\mathcal M}_{(-\infty,t)},\widehat G)$ be the maximal extension of
$({\mathcal M}_{(-\infty,t)},G)$ to time $t$, as given in Definition~\ref{hatM}.

\noindent{\bf Assumption (5). Boundary components of the exposed
regions:} {\em There is a surgery control parameter
function,
$\delta(t)>0$, a non-increasing function of
$t$, such that each component of $\Sigma_{t}\subset M_{t}$ is the
central $2$-sphere of a strong $\delta(t)$-neck in $(\widehat{\mathcal
M}_{(-\infty,t)},\widehat G)$.}


\bigskip

Suppose that $t$ is a singular time. Then for all $t^-<t$ with $t^-$
sufficiently close to $t$,  the manifolds $M_{t^-}$ are diffeomorphic and are
identified under the flow. Applying the flow (backward) to $C_{t}$ produces a
diffeomorphism from $C_{t}$  onto a compact submanifold with boundary
$C_{t^-}\subset M_{t^-}$. Our next assumption concerns the nature of the
metrics $G(t^-)$ on the {\em disappearing region}
$D_{t^-}=M_{t^-}\setminus C_{t^-}$. The following holds for
every $t^-<t$ sufficiently close to $t$.


\noindent{\bf Assumption (6). Control on the disappearing region:} {\em For any
singular time $t$, for all $t^-<t$ sufficiently close to $t$, each point of
$x\in D_{t^-}$ has a strong $(C,\epsilon)$-canonical neighborhood in
$M_{t^-}$.}


\noindent{\bf Assumption (7). Maximal flow intervals:} {\em Let $t$ be
the initial time or a singular time and let $t'$ be the first
singular time after $t$ if such exists, otherwise let $t'$ be the
least upper bound of the time-interval of definition of the Ricci
flow with surgery. Then the restriction of the Ricci flow with
surgery to $[t,t')$ is a maximal Ricci flow. That is to say, either
$t'=\infty$ or, as $t\rightarrow t'$ from below, the curvature of $G(t)$
is unbounded so that this restricted Ricci flow cannot be extended
{\bf as a Ricci flow} to any larger time.}

\bigskip

{\bf From now on $C$ and $\epsilon$ have fixed values as described
above and all Ricci flows with surgeries are implicitly assumed to
satisfy Assumptions (1) -- (7).}




\section{Topological consequences of  Assumptions (1) -- (7)}

Next we show that the topological control that we are imposing on
the $3$-dimensional Ricci flows with surgery are enough to allow
us to relate the topology of a time-slice $M_T$ in terms of a later
time-slice $M_{T'}$ and topologically standard pieces. This is the
result that will be used to establish the topological theorems
stated in the introduction.


\begin{proposition}[Topological effect of a surgery time on the earlier time-slice]
\label{prop:surgery-topological-type}\label{surgerytoptype}
Suppose that $({\mathcal M},G)$ is a generalized Ricci flow satisfying Assumptions
(1) -- (7). Let $t$ be a singular time. Then the following holds
for any $t^-<t$ sufficiently close to $t$. The manifold $M_{t^-}$ is
diffeomorphic to a manifold obtained in the following way. Take the disjoint
union of $M_t$, finitely many $2$-sphere bundles over $S^1$, and finitely many
closed $3$-manifolds admitting  metrics of constant positive curvature. Then
perform connected sum operations between (some subsets of) these components.
\end{proposition}

\begin{proof}
\uses{prop:classification-covered-by-necks-caps,prop:balanced-chain-covers-connected-set}
Fix $t'<t$ but sufficiently close to $t$.
By Assumption 4 every component of $E_t$ is a $3$-ball and hence every component
of $\partial E_t=\partial C_t$ is a $2$-sphere. Since $C_t$ is diffeomorphic to
 $C_{t'}\subset M_{t'}$ we see that every component of $\partial C_{t'}=\partial
 D_{t'}$ is a $2$-sphere. Since every component of $E_t$ is a $3$-ball,
 the passage from the smooth manifold $M_{t'}$ to the smooth manifold
 $M_t$ is effected by removing the
 interior of $D_{t'}$ from $M_{t'}$ and gluing a $3$-ball onto each component of $\partial C_{t'}$
 to form $M_t$.

By Assumption (5) every point of $D_{t'}$ has a strong
$(C,\epsilon)$-canonical neighborhood. Since $\epsilon$ is
sufficiently small it follows from Proposition~\ref{epstopology}
that every component of $D_{t'}$ that is also a component of
$M_{t'}$ is diffeomorphic either to a manifold admitting a metric of
constant positive curvature (a $3$-dimensional space-form), to
$\Ar P^3\#\Ar P^3$ or to a $2$-sphere bundle over $S^1$. In the passage from
$M_{t'}$ to $M_t$ these components are removed.

Now let us consider a component of $D_{t'}$ that is not a component
of $M_{t'}$. Such a component is a connected subset of $M_{t'}$ with
the property that every point is either contained in the core of a
$(C,\epsilon)$-cap or is the center of an $\epsilon$-neck and whose
frontier in $M_{t'}$ consists of $2$-spheres that are central
$2$-spheres of $\epsilon$-necks. If every point is the center of an
$\epsilon$-neck, then according to Proposition~\ref{epschains}
$D_{t'}$ is an $\epsilon$-tube and in particular is diffeomorphic to
$S^2\times I$. Otherwise $D_{t'}$ is contained in a capped or double
capped $\epsilon$-tube. Since the frontier of $D_{t'}$ is non-empty
and is the union of central $2$-spheres of an $\epsilon$-neck, it
follows that either $D_{t'}$ is diffeomorphic to a capped
$\epsilon$-tube or to an $\epsilon$-tube. Hence, these components of
$D_{t'}$ are diffeomorphic either to $S^2\times (0,1)$, to $D^3$, or
to $\Ar P^3\setminus B^3$. Replacing a $3$-ball component of
$D_{t'}$ by another $3$-ball leaves the topology unchanged.
Replacing a component of $D_{t'}$ that is diffeomorphic to
$S^2\times I$ by the disjoint union of two $3$-balls  has the
effect of doing a surgery along the core $2$-sphere of the cylinder
$S^2\times I$ in $M_{t'}$. If this $2$-sphere separates $M_{t'}$
into two pieces then doing this surgery effects a connected sum
decomposition. If this $2$-sphere does not separate, then the
surgery has the topological effect of doing a connected sum
decomposition into two pieces, one of which is diffeomorphic to
$S^2\times S^1$, and then removing that component entirely.
Replacing a component of $D_{t'}$ that is diffeomorphic to $\Ar
P^3\setminus B^3$ by a $3$-ball, has the effect of doing a
connected sum decomposition on $M_{t'}$ into pieces, one of which is
diffeomorphic to $\Ar P^3$, and then removing that component.

 From this description the proposition follows immediately.
\end{proof}


\begin{corollary}[Geometrization is inherited across surgery times]
\label{cor:geometrization-inherited-across-surgery}
\uses{prop:surgery-topological-type}\label{Tgood0good}
Let $({\mathcal M},G)$ be a generalized Ricci flow satisfying Assumptions
(1) -- (7)  with initial conditions $(M,g(0))$. Suppose
that for some $T$ the time-slice $M_T$ of this generalized flow
satisfies Thurston's Geometrization Conjecture. Then the same is true for the manifold
$M_t$ for any $t\le T$, and in particular $M$ satisfies Thurston's
Geometrization Conjecture. In addition,
\begin{enumerate}
\item[(1)] If for some $T>0$ the manifold $M_T$ is empty, then $M$ is a
connected sum of manifolds diffeomorphic to $2$-sphere bundles over
$S^1$ and $3$-dimensional space-forms, i.e., compact
$3$-manifolds that admit a metric of constant positive curvature.
\item[(2)]
If for some $T>0$ the manifold $M_T$ is empty and if $M$ is connected and simply connected,
then $M$ is diffeomorphic to $S^3$.
\item[(3)]
If for some $T>0$ the manifold $M_T$ is empty and if $M$ has finite fundamental group, then $M$
is a $3$-dimensional space-form.
\end{enumerate}
\end{corollary}

\begin{proof}
Suppose that $M_T$ satisfies the Thurston Geometrization Conjecture and that $t_0$ is
the largest surgery time $\le T$. (If there is no such surgery time
then $M_T$ is diffeomorphic to $M$ and the result is established.)
Let $T'<t_0$ be sufficiently close to $t_0$ so that
$t_0$ is the only surgery time in the interval $[T',T]$. Then according to
the previous proposition $M_{T'}$ is obtained from $M_T$ by first taking the disjoint union
 of $M_T$ and copies of $2$-sphere bundles over $S^1$ and $3$-dimensional space forms.
 In the Thurston Geometrization Conjecture the first step is to decompose the manifold as
 a connected sum of prime $3$-manifolds and then to treat each prime piece independently.
 Clearly, the prime decomposition of $M_{T'}$ is obtained from the prime decomposition of $M_T$
 by adding a disjoint union with $2$-sphere bundles over $S^1$ and $3$-dimensional space forms.
 By definition
 any $3$-dimensional space-form satisfies Thurston's Geometrization Conjecture.
 Since any diffeomorphism of $S^2$ to itself is isotopic to either the identity
 or to the antipodal map, there are two diffeomorphism types of $2$-sphere bundles over
 $S^1$: $S^2\times S^1$ and the non-orientable $2$-sphere bundle over $S^1$.
 Each is obtained from $S^2\times I$ be gluing the ends together by an isometry of the round
 metric on $S^2$. Hence, each has a homogeneous geometry modeled on $S^2\times \Ar$, and hence
 satisfies Thurston's Geometrization Conjecture. This proves that if $M_T$ satisfies this
 conjecture, then so does $M_{T'}$. Continuing this way by induction, using the fact that there
 are only finitely many surgery times completes the proof of the first statement.

Statement (1) is proved analogously. Suppose that $M_T$ is a disjoint union of
connected sums of $2$-sphere bundles over $S^1$
and $3$-dimensional space-forms. Let $t_0$ be
the largest surgery time $\le T$ and let $T'<t_0$ be sufficiently close to $t_0$.
(As before, if there is no such $t_0$ then $M_T$ is diffeomorphic to $M$ and the result is
established.) Then it is clear from the previous proposition that $M_{T'}$ is also a disjoint
union of connected sums of $3$-dimensional space-forms and $2$-sphere bundles over $S^1$.
Induction as in the previous case completes the argument for this case.

The last two statements are immediate from this one.
\end{proof}

\section{Further conditions on surgery}




\subsection{The surgery parameters}

The process of doing surgery requires fixing the scale $h$ at which one does
the surgery. We shall have to allow this scale $h$ to be a function of time,
decreasing sufficiently rapidly with $t$. In fact, the scale is determined by
two other functions of time which also decay to zero as time goes to infinity--
a canonical neighborhood parameter $r(t)$ determining the curvature threshold
above which we have canonical neighborhoods and the surgery control parameter
$\overline\delta(t)$ determining how close to cylinders (products of the round $2$-sphere
with an interval) the regions where we do
surgery are. In addition to these functions, in order to prove inductively that
we can do surgery we need to have a non-collapsing result. The non-collapsing
parameter $\kappa>0$ also decays to zero rapidly as time goes to infinity.
Here then are the functions that will play the crucial role in defining
the surgery process.


\begin{definition}[Canonical neighborhood and surgery control parameters; surgery scale function]
\label{def:surgery-scale-function}
\uses{thm:local-surgery,thm:deep-delta-necks-in-horn}\label{hrestricted}
We have: (i) {\em a canonical neighborhood parameter},
 $r(t)>0$, and
(ii) {\em a surgery control parameter} $\overline\delta(t)>0$. We use these
to define the surgery scale function $h(t)$. Set
$\rho(t)=\overline\delta(t)r(t)$. Let $h(t)=
h(\rho(t),\overline\delta(t))\le
\rho(t)\cdot\overline\delta(t)=\overline\delta^2(t)r(t)$ be the
function given by Theorem~\ref{hexist}. We require that $h(0)\le
R_0^{-1/2}$ where $R_0$ is the constant from
Theorem~\ref{LOCALSURGERY}.

 In addition, there is a
function $\kappa(t)>0$ called the {\em non-collapsing parameter}.
All three functions $r(t)$, $\overline\delta(t)$ and $\kappa(t)$ are
required to  be  positive,  non-increasing functions of $t$.
\end{definition}






We shall consider Ricci flows with surgery $({\mathcal M},G)$ that
satisfy Assumptions (1) -- (7) and also satisfy:

\noindent {\bf For any singular time $t$ the surgery at time $t$ is
performed with control $\overline\delta(t)$ and at scale
$h(t)=h(\rho(t),\overline\delta(t))$, where
$\rho(t)=\overline\delta(t)r(t)$, in the sense that each boundary
component of $C_t$ is the central $2$-sphere of a strong
$\overline\delta(t)$-neck centered at a point $y$ with
$R(y)=h(t)^{-2}$.}







There is quite a bit of freedom in the choice of these parameters.
But it is not complete freedom. They must decay rapidly enough as
functions of $t$. We choose to make $r(t)$ and $\kappa(t)$  step
functions, and we require $\overline\delta(t)$ to be bounded above
by a step function of $t$. Let us fix the step sizes.


\begin{definition}[Time steps T_i = 2^i t_0]
\label{def:time-steps-T-i}
We set $t_0=2^{-5}$, and for any $i\ge 0$ we define $T_i=2^it_0$.
\end{definition}

The steps we consider are $[0,T_0]$ and then $[T_i,T_{i+1}]$ for
every $i\ge 0$. The first step  is somewhat special. Suppose that
$({\mathcal M},G)$ is a Ricci flow with surgery with normalized
initial conditions. Then
 according to
Claim~\ref{newkappa0r0t0} the flow exists on $[0,T_1]$ and the norm
of the Riemann curvature is bounded by $2$ on $[0,T_1]$, so that
by Assumption (7) there are no surgeries in this time interval. Also,
by Claim~\ref{newkappa0r0t0} there is a $\kappa_0>0$ so
that $\Vol\,B(x,t,r)\le \kappa_0r^3$ for every $t\le T_1$ and
$x\in M_t$ and every $r\le \epsilon$.

\begin{definition}[Surgery parameter sequences (r, kappa, delta)]
\label{def:surgery-parameter-sequences}
\uses{thm:local-surgery,prop:neck-gluing,claim:normalized-initial-existence,lem:standard-initial-metric-asymptotic-bounds}\label{surgeryparams}
{\em Surgery parameter sequences} are sequences
\begin{enumerate}
\item[(i)] ${\bf r}=r_0\ge r_1\ge r_2\ge \cdots >0$, with $r_0=\epsilon$,
\item[(ii)] ${\bf K}=\kappa_0\ge \kappa_1\ge \kappa_2\ge \cdots >0$
with $\kappa_0$ as in Claim~\ref{newkappa0r0t0}, and \item[(iii)] $
\Delta=\delta_0\ge \delta_1\ge \delta_2\ge \cdots
>0$ with $\delta_0=\mathrm{min}(\beta\epsilon/3,\delta_0',K^{-1},D^{-1})$
 where $\delta_0'$ is the constant
from Theorem~\ref{LOCALSURGERY} and $\beta<1/2$ is the constant from
Proposition~\ref{neckglue}, $\epsilon$ is the constant that we have
already fixed, and $K$ and $D$ are the constants from
Lemma~\ref{A_0}.
\end{enumerate}
We shall also refer to partial sequences defined for indices $0,\ldots, i$ for some $i>0$
as surgery parameter sequences if they are positive, non-increasing and if their initial
terms satisfy the conditions given above.
\end{definition}


We let $r(t)$ be the step function whose value on $[T_i,T_{i+1})$ is $r_{i+1}$
and whose value on $[0,T_0)$ is $r_0$. We say that a Ricci flow with surgery
satisfies the strong $(C,\epsilon)$-canonical neighborhood assumption with parameter ${\bf
r}$ if it satisfies this condition with respect to the step function $r(t)$
associated with ${\bf r}$. This means that any $x\in {\mathcal M}$ with
$R(x)\ge r^{-2}({\bf t}(x))$ has a strong $(C,\epsilon)$-canonical neighborhood in
${\mathcal M}$. Let $\kappa(t)$ be the step function whose value on
$[T_i,T_{i+1})$ is $\kappa_{i+1}$ and whose value on $[0,T_0)$ is $\kappa_0$.
Given $\kappa>0$, we say that a Ricci flow defined on $[0,t]$ is $\kappa$-non-collapsed
on scales $\le \epsilon$ provided that for every point $x$ not contained in a
component of its time-slice with positive sectional curvature,
if for some $r\le \epsilon$, the parabolic neighborhood $P(x,{\bf t}(x),r,-r^2)$ exists in ${\mathcal M}$
and the norm of the Riemann curvature is bounded on this backward
parabolic neighborhood by $r^{-2}$, then $\Vol\,B(x,{\bf t}(x),r)\ge
\kappa r^3$.
We say that a Ricci flow with surgery is ${\bf K}$-non-collapsed on scales
$\epsilon$ if for every $t\in [0,\infty)$ the restriction of the flow to
$[0,t]$ is $\kappa(t)$-non-collapsed on scales $\le \epsilon$.
  Lastly, we fix a non-increasing function $\overline\delta(t)>0$
with $\overline\delta(t)\le \delta_{i+1}$ if $t\in [T_i,T_{i+1})$ for all $i\ge
0$ and $\overline\delta(t)\le \delta_0$ for $t\in [0,T_0)$. We denote the fact
that such inequalities hold for all $t$ by saying $\overline\delta(t)\le
\Delta$.


Having fixed surgery parameter sequences ${\bf K}$, ${\bf r}$ and
$\Delta$, defined step functions $r(t)$ and $\kappa(t)$, and fixed
$\overline\delta(t)\le \Delta$ as above, we shall consider only
Ricci flows with surgery  where the surgery at time $t$ is defined
using the surgery parameter functions $r(t)$ and
$\overline\delta(t)$. In addition,  we require that these Ricci
flows with surgery satisfy Assumptions (1) -- (7).

What we shall  show is that there are surgery parameter
sequences ${\bf r}$, ${\bf K}$ and $\Delta$ with the property that
for any normalized initial metric and any positive,
non-increasing function $\overline\delta(t)\le \Delta$, it is
possible to construct a Ricci flow with surgery using the surgery
parameters $r(t)$ and $\overline\delta(t)$ with the given initial
conditions and furthermore that this Ricci flow with surgery
satisfies the Assumptions (1) -- (7), has curvature pinched toward
positive, satisfies the canonical neighborhood assumption, and satisfies the
non-collapsing assumption using these parameters.

In fact we shall prove this inductively, constructing the step
functions inductively one step at a time. Thus, given surgery
parameter sequences indexed by $0,\ldots,i$ we show that there are appropriate choices of
$r_{i+1}, \kappa_{i+1}$ and $\delta_{i+1}$ such that the following is true.
Given a Ricci flow with
surgery defined on time $[0,T_i)$ satisfying all the properties with
respect to the first set of data,  that Ricci
flow with surgery extends to one defined for time $[0,T_{i+1})$ and satisfies
Assumptions (1) -- (7), the canonical neighborhood assumption and the
non-collapsing assumption with respect to the extended surgery
parameter sequences, and has curvature pinched toward positive.
As stated this is not quite true; there is a slight twist: we must also assume
that $\overline\delta(t)\le \delta_{i+1}$ for all $t\in [T_{i-1},T_{i+1})$.
It is for this reason that we consider
pairs consisting of sequences $\Delta$ and a surgery control
parameter $\overline\delta(t)$ bounded above by $\Delta$.


\section{The process of surgery}\label{sect:process}


Suppose given surgery parameter sequences $\{r_0,\ldots,r_{i+1}\}$,
$\{\kappa_0,\ldots,\kappa_{i+1}\}$ and
$\Delta_i=\{\delta_0,\ldots,\delta_i\}$ and also given a positive,
decreasing function $\overline\delta(t)\le \Delta_i$, defined for
$t\le T_{i+1}$ with $\delta_0=\mathrm{min}(\alpha\epsilon/3,\delta_0',K^{-1},D^{-1})$ as above. Suppose
that $({\mathcal M},G)$ is a Ricci flow with surgery defined for
$t\in [0,T)$ that goes singular at time $T\in (T_i,T_{i+1}]$. We
suppose that it satisfies Assumptions (1) -- (7). Since the flow has
normalized initial conditions and goes singular at time $T$, it
follows that $i\ge 1$. We suppose that $({\mathcal M},G)$ satisfies
the $(C,\epsilon)$-canonical neighborhood assumption with parameter
$r_{i+1}$ and that its curvature is pinched toward positive. By
Theorem~\ref{omega} we know that there is a maximal extension
$(\widehat{\mathcal M},\widehat G)$ of this generalized flow to time
$T$ with the $T$ time-slice being $(\Omega(T),G(T))$. Set
$\rho=\overline\delta(T)r_{i+1}$, and set
$h(T)=h(\rho(T),\overline\delta(T))$ as in Theorem~\ref{hexist}.
Since $\overline\delta(T)\le \delta_0<1$, we see that
$\rho<r_{i+1}$. According to Lemma~\ref{2epshorns} there are
finitely many components of $\Omega(T)$ that meet $\Omega_\rho(T)$.
Let $\Omega^\mathrm{big}(T)$ be the disjoint union of all the
components of $\Omega(T)$ that meet $\Omega_\rho(T)$.
Lemma~\ref{2epshorns} also tells us that $\Omega^\mathrm{big}(T)$
contains a finite collection of disjoint $2\epsilon$-horns with
boundary contained in $\Omega_{\rho/2C}$, and the complement of the
union of the interiors of these horns is a compact submanifold with
boundary containing $\Omega_\rho$.  Let ${\mathcal
H}_1,\ldots,{\mathcal H}_j$ be a disjoint union of these
$2\epsilon$-horns. For each $i$ fix a point $y_i\in {\mathcal H}_i$
with $R(y_i)=h^{-2}(T)$. According to Theorem~\ref{hexist} for each
$i$ there is a strong $\delta(T)$-neck centered at $y_i$ and
contained in ${\mathcal H}_i$.
 We orient the $s$-direction of the
neck so that its positive end lies closer to the end of the horn
than its negative end. Let $S^2_i$ be the center of this strong
$\delta(T)$-neck. Let ${\mathcal H}_i^+$ be the unbounded
complementary component of $S^2_i$ in ${\mathcal H}_i$. Let $C_T$ be
the complement of $\coprod_{i=1}^j{\mathcal H}^+_i$ in $\Omega^\mathrm{big}(T)$. Then we do surgery on these necks as described in
Section~\ref{necksurgery}, using the constant $q=q_0$ from
Theorem~\ref{LOCALSURGERY},
 removing the positive half of the neck,
and gluing on the cap from the standard solution. This creates a compact
$3$-manifold $M_T=C_T\cup_{\coprod_iS^2_i} B_i$, where each $B_i$ is a copy of the
metric ball of radius $A_0+4$ centered around the tip of the standard solution
 (with the metric scaled by
$h^2(T)$ and then perturbed near the boundary of $B_i$ to
%
match $g(T)$). Notice that in this process we have removed every
component of $\Omega(T)$ that does not contain a point of
$\Omega_\rho(T)$. The result of this operation is to produce a
compact Riemannian $3$-manifold $(M_T,G_T)$ which is the $T$
time-slice of our extension of $({\mathcal M},G)$. Let $(M_T,G(t)),\
T\le t<T'$, be the maximal Ricci flow with initial conditions
$(M_T,G_T)$ at $t=T$. Our new space-time is the union of
$M_T\times [T,T')$ and $({\mathcal M},G)\cup C_T\times \{T\}$  along $C_T\times \{T\}$.
Since the isometric embedding  $C_T\subset M_T$ extends to an
Here, we view $({\mathcal M},G)\cup C_T\times \{T\}$ as a subspace
of $(\widehat{\mathcal M},\widehat G)$ via the
isometric embedding of $C_T$ into $\Omega(T)$. The time functions
and vector fields glue to provide analogous data for this new
space-time. According to Lemma~\ref{gluingsoln} the horizontal
metrics glue together to make a smooth metric on space-time
satisfying the Ricci flow equation.


Notice that the continuing region at time $T$ is exactly $C_T$
whereas the exposed region is $\coprod_iB_i$, which is a disjoint
union of $3$-balls. The disappearing region is the complement of
the embedding of $C_T$ in $M_{t'}$ for $t'<T$ but sufficiently close to it obtained by
flowing $C_T\subset \Omega_T$ backward. The
disappearing region contains $M_{t'}\setminus \Omega(T)$ and also
contains all components of $\Omega(T)$ that do not contains points
of $\Omega_\rho(T)$, as well as the ends of those components of
$\Omega(T)$ that contain points of $\Omega_\rho(T)$.

\begin{definition}[The surgery operation at a singular time]
\label{def:surgery-operation}
\uses{thm:local-surgery,lem:gluing-ricci-flows-with-surgery,thm:singular-time-limit-properties,thm:deep-delta-necks-in-horn,lem:finitely-many-horns-meeting-omega-rho}
The operation described in the previous paragraph is the {\em surgery
operation at time $T$} using the surgery parameters $\overline\delta(T)$ and $r_{i+1}$.
\end{definition}


\section{Statements about the existence of Ricci flow with
surgery}

What we shall establish is the existence of surgery satisfying Assumptions (1) -- (7)
above and also satisfying the curvature pinched toward positive assumption, the
strong canonical neighborhood assumption, and the non-collapsing assumption.
This requires first of all that we begin with a compact, Riemannian
$3$-manifold $(M,g(0))$ that is normalized, which we are assuming. It also
requires careful choice of upper bounds $\Delta=\{\delta_i\}$ for the surgery
control parameter $\overline\delta(t)$ and careful choice of the canonical neighborhood
parameter ${\bf r}=\{r_i\}$ and of the non-collapsing step function ${\bf
K}=\{\kappa_i\}$.

Here is the statement that we shall establish.


\begin{theorem}[Main existence theorem: controlled Ricci flow with surgery extends for all time]
\label{thm:main-existence-ricci-flow-surgery}\label{MAIN}
\uses{def:curvature-pinched-toward-positive,def:strong-canonical-neighborhood-assumption}
There are surgery parameter sequences
$${\bf
K}=\{\kappa_i\}_{i=1}^\infty,\Delta=\{\delta_i\}_{i=1}^\infty,{\bf
r}=\{r_i\}_{i=1}^\infty$$
  such that the following holds. Let $r(t)$
be the step function whose value on $[T_{i-1},T_i)$ is $r_i$.
Suppose that $\overline\delta\colon [0,\infty)\to \Ar^+$ is any
non-increasing function with $\overline\delta(t)\le \delta_i$
whenever $t\in [T_{i-1},T_i)$. Then the following holds: Suppose
that $({\mathcal M},G)$ is a Ricci flow with surgery defined for
$0\le t<T$ satisfying Assumptions (1) -- (7). In addition, suppose that
the following conditions:
\begin{enumerate}
\item[(1)] the generalized flow has curvature
pinched toward positive,
\item[(2)] the flow satisfies the strong $(C,\epsilon)$-canonical neighborhood assumption
with parameter ${\bf r}$ on $[0,T)$, and \item[(3)] the flow is ${\bf K}$
non-collapsed on $[0,T)$ on scales $\le \epsilon$.
\end{enumerate}
 Then there is
an extension of $({\mathcal M},G)$ to a Ricci flow with surgery
defined for all $0\le t<\infty$ and satisfying Assumptions (1) -- (7) and
the above three conditions.
\end{theorem}

This of course leads immediately to the existence result for Ricci
flows with surgery defined for all time with any normalized initial
conditions.


\begin{corollary}[Every normalized compact 3-manifold admits a Ricci flow with surgery for all time]
\label{cor:ricci-flow-surgery-exists-all-time}\label{RFSexists}
Let ${\bf K}$, ${\bf r}$ and $\Delta$ be surgery parameter sequences
provided by the previous theorem. Let $\overline\delta(t)$ be a
non-increasing positive function with $\overline\delta(t)\le
\Delta$. Let $M$ be a compact $3$-manifold containing no $\Ar P^2$
with trivial normal bundle. Then there is a Riemannian metric $g(0)$
on $M$ and a Ricci flow with surgery defined for $0\le t<\infty$
with  initial metric $(M,g(0))$. This Ricci flow with surgery
satisfies the seven assumptions and is ${\bf K}$-non-collapsed on
scales $\le \epsilon$. It also satisfies the strong
$(C,\epsilon)$-canonical neighborhood assumption with parameter
${\bf r}$ and has curvature pinched toward positive. Furthermore,
any surgery at a time $t\in [T_i,T_{i+1})$ is done using
$\overline\delta(t)$ and $r_{i+1}$.
\end{corollary}

\begin{proof}
\uses{thm:main-existence-ricci-flow-surgery,prop:normalized-flow-short-time-bounds} (Assuming Theorem~\ref{MAIN})
Choose a metric $g(0)$ so that $(M,g_0)$ is normalized. This is
possible by beginning with any Riemannian metric on $M$ and scaling
it by a sufficiently large positive constant to make it normalized.
According to Proposition~\ref{kappa0r0t0} and the definitions of
$T_i$ and $\kappa_0$ there is a Ricci flow $(M,g(t))$ with these
initial conditions defined for $0\le t\le T_2$ satisfying Assumptions (1) -- (7)
and the three conditions of the previous theorem. The assumption that
$M$ has no embedded $\Ar P^2$ with trivial normal bundle is needed so that Assumption (1) holds
for this Ricci flow.
 Hence, by the previous theorem we can extend this
Ricci flow to a Ricci flow with surgery defined for all $0\le
t<\infty$ satisfying the same conditions.
\end{proof}





Showing that after surgery Assumptions (1) -- (7) continue to hold and
that the curvature is pinched toward positive is direct and only requires that
$\overline\delta(t)$ be smaller than some universal positive constant.

\begin{lemma}[Surgery preserves Assumptions (1)-(7) and curvature pinching]
\label{lem:surgery-preserves-assumptions-pinching}\label{14.26}
Suppose that $({\mathcal M},G)$ is a Ricci flow with surgery going
singular at time $T\in [T_{i-1},T_i)$. We suppose that $({\mathcal
M},G)$ satisfies Assumptions (1) - (7),  has curvature pinched toward
positive, satisfies the strong $(C,\epsilon)$-canonical neighborhood
assumption with parameter ${\bf r}$ and is ${\bf K}$ non-collapsed.
Then the result of the surgery operation  at time $T$ on $({\mathcal M},G)$ is a
 Ricci flow with surgery defined on $[0,T')$ for some $T'>T$.
 The resulting Ricci flow with
surgery satisfies Assumptions (1) -- (7). It also has curvature pinched
toward positive.
\end{lemma}

\begin{proof}
\uses{thm:local-surgery,lem:gluing-ricci-flows-with-surgery,thm:pinching-persists-from-time-a}
It is immediate from the construction and  Lemma~\ref{gluingsoln} that
the result of performing the surgery operation at time $T$ on a Ricci flow with
surgery produces a new Ricci flow with surgery. Assumptions (1) -- (3)
clearly hold for the result. and Assumptions (4) and (5) hold because of
the way that we do surgery. Let us consider Assumption (6). Fix $t'<T$
so that there are no surgery times in $[t',T)$. By flowing backward
using the vector field $\chi$ we have an embedding $\psi\colon
C_t\times [t',T]\to \widehat {\mathcal M}$ compatible with time and
the vector field. For any $p\in M_{t'}\setminus\psi(\mathrm{int}\,C_T\times \{t'\})$ the limit as $t$ tends to $T$ from below of
the flow line $p(t)$ at time $t$ through $p$ either lies in
$\Omega(T)$ or it does not. In the latter case, by definition we
have
$$\mathrm{lim}_{t\rightarrow T^-}R(p(t))=\infty.$$
In the former case, the limit point  either is contained in the end
of a strong $2\epsilon$-horn cut off by the central $2$-sphere of
the strong $\delta$-neck centered at one of the $y_i$ or is
contained in a component of $\Omega(T)$ that contains no point of
$\Omega_\rho(T)$. Hence, in this case we have
$$\mathrm{lim}_{t\rightarrow T^-}R(p(t))>\rho^{-2}>r_i^{-2}.$$
Since $M_{t'}\setminus \psi(\mathrm{int}\,C_T\times \{t'\})$ is compact
for every $t'$, there is $T_1<T$ such that $R(p(t))>r_i^{-2}$ for
all $p\in M_{t'}\setminus \psi(\mathrm{int}C_T\times \{t'\})$ and all
$t\in [T_1,T)$. Hence, by our assumptions all these points have
strong $(C,\epsilon)$-canonical neighborhoods. This establishes that
Assumption (6) holds at the singular time $T$. By hypothesis
Assumption (6) holds at all earlier singular times. Clearly, from the
construction the Ricci flow on $[T,T')$ is maximal. Hence,
Assumption (7) holds for the new Ricci flow with surgery.

 From Theorem~\ref{LOCALSURGERY} the
fact that $\delta(T)\le \delta_i\le \delta_0\le\delta'_0$ and
$h(T)\le R_0^{-1/2}$ imply that the Riemannian manifold $(M_T,G(T))$
has curvature pinched toward positive for time $T$. It then follows
from Corollary~\ref{pincha} that  the Ricci flow defined on $[T,T')$
with $(M_T,G(T))$ as initial conditions has curvature pinched toward
positive. The inductive hypothesis is that on the time-interval
$[0,T)$ the Ricci flow with surgery has curvature pinched toward
positive. This completes the proof of the lemma.
\end{proof}

\begin{proposition}[The pre-surgery flow map is distance decreasing near a surgery time]
\label{prop:surgery-distance-decreasing}
\uses{thm:local-surgery,def:surgery-parameter-sequences}\label{surgerydist}
Suppose that $({\mathcal M},G)$ is a Ricci flow with surgery
satisfying Assumptions (1) -- (7) in Section~\ref{firstassumpt}.
Suppose that $T$ is a surgery time, suppose  that the surgery
control parameter $\delta(T)$ is less than $\delta_0$ in
Definition~\ref{surgeryparams}, and suppose that the scale of the
surgery $h(T)$ is less than $R_0^{-1/2}$ where $R_0$ is the constant
from Theorem~\ref{LOCALSURGERY}. Fix $t'<T$  sufficiently close to
$T$. Then  there is an embedding $\rho\colon M_{t'}\times [t',T)\to
{\mathcal M}$ compatible with time and the vector field. Let $X(t')$
be a component of $M_{t'}$ and let $X(T)$ be a component obtained
from $X(t')$ by doing surgery at time $T$. We view $\rho^*G$ as a
one-parameter family of metrics $g(t)$ on $X(t')$. There is an open
subset $\Omega\subset X(t')$ with the property that $\mathrm{lim}_{t'\rightarrow T^-} g(t')|_\Omega$ exists (we denote it by
$g(T)|_\Omega$) and with the property that $\rho|_{\Omega\times
[t',T)}$ extends to a map $\widehat \rho\colon \Omega\times
[t',T]\to {\mathcal M}$. This defines a map for $\Omega\subset
X(t')$ onto an open subset $\Omega(T)$ of $X(T)$ which is an
isometry from the limiting metric $g(T)$ on $\Omega$ to
$G(T)|_{\Omega}$. Suppose that all of the $2$-spheres along which we
do surgery are separating. Then this map extends to a map $X(t')\to
X(T)$. For all $t<T$ but sufficiently close to $T$ this extension is
a distance decreasing map from $(X(t')\setminus \Omega,g(t))$ to
$X(T)$.
\end{proposition}

\begin{proof}
\uses{thm:local-surgery}
This is immediate from the third item in Theorem~\ref{LOCALSURGERY}.
\end{proof}


\begin{remark}[Non-separating surgery spheres prevent extending the comparison map]
\label{rem:nonseparating-surgery-sphere}
\uses{prop:surgery-distance-decreasing}
If we have a non-separating surgery $2$-sphere then there will a
component $X(T)$ with surgery caps on both sides of the surgery
$2$-sphere and hence we cannot extend the map even continuously
over all of $X(t')$.
\end{remark}

The other two inductive properties in Theorem~\ref{MAIN} -- that the result is
${\bf K}$-non-collapsed and also that it satisfies the strong
$(C,\epsilon)$-canonical neighborhood assumption with parameter ${\bf r}$ -
require appropriate inductive choices of the sequences. The arguments
establishing these are quite delicate and intricate. They are given in the next
two sections.

\section{Outline of the proof of Theorem~{\ref{MAIN}}}

Before giving the proof proper of Theorem~\ref{MAIN} let us outline
how the argument goes. We shall construct the surgery parameter
sequences $\Delta$, ${\bf r}$, and ${\bf K}$ inductively. Because of
Lemma~\ref{kappa0r0t0} we have the beginning of the inductive
process. We suppose that we have defined sequences as required up to
index $i$ for some $i\ge 1$. Then we shall extend them one more step
to sequences defined up to $(i+1)$, though there is a twist: to do
this we must redefine $\delta_i$ in order to make sure that the
extension is possible. In Chapter~\ref{sectnoncoll} we establish the
non-collapsing result assuming the strong canonical neighborhood
result. More precisely, suppose that we have a Ricci flow with
surgery $({\mathcal M},G)$ defined for time $0\le t<T$ with $T\in
(T_i,T_{i+1}]$ so that the restriction of this flow to the
time-interval $[0,T_i)$ satisfies the inductive hypothesis with
respect to the given sequences. Suppose also that the entire Ricci
flow with surgery has
 strong $(C,\epsilon)$-canonical neighborhoods for some $r_{i+1}>0$.
Then there is $\delta(r_{i+1})>0$ and $\kappa_{i+1}>0$ such that,
provided that $\overline\delta(t)\le \delta(r_{i+1})$ for all $t\in
[T_{i-1},T)$, the Ricci flow with surgery $({\mathcal M},G)$ is
$\kappa_{i+1}$ non-collapsed on scales $\le \epsilon$.

In Section~\ref{sectcannbhd} we show that the strong $(C,\epsilon)$-canonical neighborhood
assumption extends for some parameter $r_{i+1}$, assuming again that
$\overline\delta(t)\le \delta(r_{i+1})$ for all $t\in [T_{i-1},T)$.

Lastly, in Section~\ref{sectcomplete} we complete the proof by
showing that the number of surgeries possible in $[0,T_{i+1})$ is
bounded in terms of the initial conditions and $\overline\delta(T)$.
The argument for this is a simple volume comparison argument --
under Ricci flow with normalized initial conditions, the volume
grows at most at a fixed exponential rate and under each surgery an
amount of volume, bounded below by a positive constant depending
only on $\overline \delta(T_{i+1})$, is removed.

